\documentclass[10pt,a4paper]{amsart}
\usepackage{amsmath,amssymb,amsthm,graphicx}
\usepackage{verbatim}
\newenvironment{reference}{\comment}{\endcomment}
%\newenvironment{reference}{}{}
\newenvironment{slogan}{\comment}{\endcomment}
\newenvironment{history}{\comment}{\endcomment}

% For commutative diagrams we use Xy-pic
\usepackage[all]{xy}

% We use 2cell for 2-commutative diagrams.
\xyoption{2cell}
% Theorem environments.
%
\theoremstyle{plain}
\newtheorem{theorem}[subsection]{Theorem}
\newtheorem{proposition}[subsection]{Proposition}
\newtheorem{lemma}[subsection]{Lemma}

\theoremstyle{definition}
\newtheorem{definition}[subsection]{Definition}
\newtheorem{example}[subsection]{Example}
\newtheorem{exercise}[subsection]{Exercise}
\newtheorem{situation}[subsection]{Situation}

\theoremstyle{remark}
\newtheorem{remark}[subsection]{Remark}
\newtheorem{remarks}[subsection]{Remarks}

\numberwithin{equation}{subsection}

% Macros
%
\def\lim{\mathop{\mathrm{lim}}\nolimits}
\def\colim{\mathop{\mathrm{colim}}\nolimits}
\def\Spec{\mathop{\mathrm{Spec}}}
\def\Hom{\mathop{\mathrm{Hom}}\nolimits}
\def\Ext{\mathop{\mathrm{Ext}}\nolimits}
\def\SheafHom{\mathop{\mathcal{H}\!\mathit{om}}\nolimits}
\def\SheafExt{\mathop{\mathcal{E}\!\mathit{xt}}\nolimits}
\def\Sch{\mathit{Sch}}
\def\Mor{\mathop{\mathrm{Mor}}\nolimits}
\def\Ob{\mathop{\mathrm{Ob}}\nolimits}
\def\Sh{\mathop{\mathit{Sh}}\nolimits}
\def\NL{\mathop{N\!L}\nolimits}
\def\CH{\mathop{\mathrm{CH}}\nolimits}
\def\proetale{{pro\text{-}\acute{e}tale}}
\def\etale{{\acute{e}tale}}
\def\QCoh{\mathit{QCoh}}
\def\Ker{\mathop{\mathrm{Ker}}}
\def\Im{\mathop{\mathrm{Im}}}
\def\Coker{\mathop{\mathrm{Coker}}}
\def\Coim{\mathop{\mathrm{Coim}}}

% Boxtimes
%
\DeclareMathSymbol{\boxtimes}{\mathbin}{AMSa}{"02}

%
% Macros for moduli stacks/spaces
%
\def\QCohstack{\mathcal{QC}\!\mathit{oh}}
\def\Cohstack{\mathcal{C}\!\mathit{oh}}
\def\Spacesstack{\mathcal{S}\!\mathit{paces}}
\def\Quotfunctor{\mathrm{Quot}}
\def\Hilbfunctor{\mathrm{Hilb}}
\def\Curvesstack{\mathcal{C}\!\mathit{urves}}
\def\Polarizedstack{\mathcal{P}\!\mathit{olarized}}
\def\Complexesstack{\mathcal{C}\!\mathit{omplexes}}
% \Pic is the operator that assigns to X its picard group, usage \Pic(X)
% \Picardstack_{X/B} denotes the Picard stack of X over B
% \Picardfunctor_{X/B} denotes the Picard functor of X over B
\def\Pic{\mathop{\mathrm{Pic}}\nolimits}
\def\Picardstack{\mathcal{P}\!\mathit{ic}}
\def\Picardfunctor{\mathrm{Pic}}
\def\Deformationcategory{\mathcal{D}\!\mathit{ef}}
\usepackage{cleveref}
\crefname{originallemma}{original lemma}{original lemmas}
\Crefname{originallemma}{Original lemma}{Original lemmas}
\makeatletter
\newlabel{lemma-localize-NL}{{00S7}{source}{}{sourcepointer-1}{}}
\newlabel{lemma-localize-NL@cref}{{[originallemma][0][]00S7}{[1][0][]source}{}{}{}}
\makeatother
\crefname{originallemma}{original lemma}{original lemmas}
\Crefname{originallemma}{Original lemma}{Original lemmas}
\makeatletter
\newlabel{lemma-relative-global-complete-intersection-conormal}{{00SV}{source}{}{sourcepointer-2}{}}
\newlabel{lemma-relative-global-complete-intersection-conormal@cref}{{[originallemma][0][]00SV}{[1][0][]source}{}{}{}}
\makeatother
\title{Smoothness of a relative global complete intersection is detected by a nonvanishing maximal minor}
\author{The Stacks Project Authors}
\date{stacks-project:algebra.tex:local-snapshot-20261012; GNU FDL 1.2 or later}
\begin{document}
\maketitle
\noindent\textit{Editorial note.} All delivered mathematical text is the author's own, copied from the pinned Stacks Project chapter file algebra.tex; no statement or proof step was rewritten, reordered or removed.\par\medskip
\noindent\textit{Editorial note.} Preamble provenance, as required for this source: the frozen snapshot stages the chapter's own preamble file beside it (preamble.tex, 217 lines, 8116 bytes, the same shared Stacks-wide preamble the published chow and etale-cohomology cards use); the first line of the pinned chapter file is \textbackslash{}input\{preamble\}, so no preamble text lives inside the chapter itself. The card takes that file's lines 8-18 (the verbatim package and the three \textbackslash{}comment-based environments reference/slogan/history at lines 9-12, section 13.2, together with the Xy-pic setup) and its lines 154 to the runtime end of the file (the native theorem declarations and mathematical macros). Not one character was changed, no package was added, no macro was defined, and the excerpt serves only as preamble, never as a span coordinate.\par\medskip
\noindent\textit{Editorial note.} The same-chapter pointers are located at their \textbackslash{}label lines inside the pinned chapter file; the cross-chapter pointers are located in the project's own tag list staged in the snapshot (sources/stacks-chow/tags.txt), at the verbatim lines carrying their canonical tags.\par\medskip
\noindent\textit{Editorial note.} The counted proof refers to two same-chapter lemmas by label; both are recorded as navigation stubs pointed at their \textbackslash{}label lines, and the arguments consuming them stay inside the counted span.\par\medskip
\noindent\textit{Editorial note.} The counted spans contain no \textbackslash{}cite, so the card carries no bibliography.\par\medskip
\noindent\textit{Allowed standard prerequisites (editorial).}\begin{itemize}
\item The shape of the naive cotangent complex of a relative global complete intersection. Invoked in the counted proof at lines 37860-37867. Identifies the complex with an explicit two-term complex given by the Jacobian matrix.
\item Localising the naive cotangent complex. Invoked twice in the counted proof (lines 37876 and 37885). Relates smoothness after inverting an element to the localized complex being split.
\end{itemize}\medskip
\ifdefined\hypertarget\hypertarget{sourcepointer-1}{}\fi
\noindent\textit{Source reference.} Original lemma 00S7: sources/stacks-algebra/pinned-original-source.tex, source lines 35909--35909. Navigation pointer for \textbackslash{}ref\{lemma-localize-NL\} used inside the counted span (line 37876); the argument that consumes it is retained in the span.\par\medskip
\ifdefined\hypertarget\hypertarget{sourcepointer-2}{}\fi
\noindent\textit{Source reference.} Original lemma 00SV: sources/stacks-algebra/pinned-original-source.tex, source lines 36984--36984. Navigation pointer for \textbackslash{}ref\{lemma-relative-global-complete-intersection-conormal\} used inside the counted span (line 37860); the argument that consumes it is retained in the span.\par\medskip
\begin{lemma}
\label{lemma-relative-global-complete-intersection-smooth}
Let $R$ be a ring. Let $S = R[x_1, \ldots, x_n]/(f_1, \ldots, f_c)$
be a relative global complete intersection.
Let $\mathfrak q \subset S$ be a prime. Then $R \to S$
is smooth at $\mathfrak q$ if and only if there exists a
subset $I \subset \{1, \ldots, n\}$ of cardinality $c$
such that the polynomial
$$
g_I = \det (\partial f_j/\partial x_i)_{j = 1, \ldots, c, \ i \in I}.
$$
does not map to an element of $\mathfrak q$.
\end{lemma}

\begin{proof}
By Lemma \ref{lemma-relative-global-complete-intersection-conormal}
we see that the naive cotangent complex
associated to the given presentation of $S$ is the complex
$$
\bigoplus\nolimits_{j = 1}^c S \cdot f_j
\longrightarrow
\bigoplus\nolimits_{i = 1}^n S \cdot \text{d}x_i, \quad
f_j \longmapsto \sum \frac{\partial f_j}{\partial x_i} \text{d}x_i.
$$
The maximal minors of the matrix giving the map are exactly
the polynomials $g_I$.

\medskip\noindent
Assume $g_I$ maps to $g \in S$, with $g \not \in \mathfrak q$.
Then the algebra $S_g$ is smooth over $R$. Namely, its naive
cotangent complex is quasi-isomorphic to the complex above
localized at $g$, see Lemma \ref{lemma-localize-NL}. And by
construction it is quasi-isomorphic to a free rank $n - c$
module in degree $0$.

\medskip\noindent
Conversely, suppose that all $g_I$ end up in $\mathfrak q$.
In this case the complex above tensored with $\kappa(\mathfrak q)$
does not have maximal rank, and hence there is no localization
by an element $g \in S$, $g \not \in \mathfrak q$
where this map becomes a split injection. By Lemma \ref{lemma-localize-NL}
again there is no such localization which is smooth over $R$.
\end{proof}

\end{document}
